% STAGING DRAFT for Direction B -- NOT yet \input into main.tex.
% Slots marked <<FILL:...>> await the full break_results.json (cross-config correlation + sharpening).
% Confirmed single-config numbers (argmax, NoPE, baseline, seed 0) are inline and final.
% Intended placement: a new section right after "Why attention concentration and variance separate"
% (the theory section), before "Restoring concentration".

\section{Predicting the break point}
\label{sec:break}

Proposition 2 says attention on the correct source holds only if the attention logit gap grows about like the
log of the context.
This gives a prediction, not only a description.
At the answer query the softmax over keys is, exactly,
\begin{equation}
a_{j^*} = \sigma\!\left(\Delta(n)\right),
\qquad
\Delta(n) = z_{j^*} - \operatorname{logsumexp}_{k \neq j^*} z_k,
\label{eq:delta}
\end{equation}
where $z$ are the effective attention logits at the retrieval head and $\sigma$ is the logistic function.
The background term grows about like $\log n$ as the context lengthens, so $a_{j^*}$ stays high only while
$z_{j^*}$ keeps pace, and the model breaks where $\Delta(n)$ reaches zero.
We fit $\Delta(n) = \alpha + \beta \log n$ on short lengths alone, inside and just past the training regime,
and read off the predicted break as the length where the fit crosses zero.
The observed break is the length where exact-match crosses one half.

The gap predicts the break (Figure~\ref{fig:break}).
On argmax retrieval without positional encoding, the gap measured only at lengths five to twenty, where
exact-match is at least $0.98$, extrapolates to a predicted break at length $65$, and the model breaks at
length $84$, within a quarter of the prediction.
The gap declines with the log of the context at a slope of $-2.4$, and the per-key logit margin does not grow
fast enough to hold it up, which is the mechanism of the failure.
Across the <<FILL:N>> configurations that break, the predicted and observed break lengths agree, with a
correlation of <<FILL:corr>> in the log length and a median ratio of <<FILL:ratio>>.
The break point is therefore forecast from a statistic measured while the model still generalizes.

The prediction also explains the intervention.
Length-scaled logit sharpening is designed to add the missing $\log n$ growth to the gap.
With sharpening the fitted slope moves from <<FILL:slope_base>> to <<FILL:slope_sharp>>, so the gap holds up
and the predicted break moves out to <<FILL:break_sharp>>, which matches the later length at which the
sharpened model fails.
The same statistic that predicts the break also predicts which intervention moves it.

\begin{figure}[H]
\centering
\includegraphics[width=\columnwidth]{figures/fig_break.pdf}
\caption{The attention logit gap predicts the length-generalization break point.
Left: the break length predicted from a short-length fit of the gap against the break length observed from
accuracy, one point per configuration, on a log-log axis.
Right: the gap $\Delta$ against the log of the context for a baseline and its sharpened twin; the baseline
crosses zero and breaks, and sharpening holds the gap up.}
\label{fig:break}
\end{figure}

% Honest caveat to fold into Limitations when this section is adopted:
% The prediction extrapolates a short-length linear-in-log-n fit, so a gap that is not linear in log n moves
% the point prediction; the break is read at a_{j*}=0.5 and exact-match 0.5, both conventions; the retrieval
% head is chosen at the shortest length. These are stated where the section is integrated.
